\section{Per-Fold Cross-Validation Results}
\label{app:perfold}

Tables~\ref{tab:perfold_v7_nll} and~\ref{tab:perfold_v8_training} report per-fold results for the V7 and V8 decoders, respectively, documenting the fold-to-fold variability underlying the mean AUROC values reported in the main text.

\begin{table}[ht]
\centering
\caption{V7 decoder per-fold AUROC for selected attack types and scoring methods. Fold-to-fold variation is modest for A1 (range = 0.154 for $S_{\text{mean}}$) but substantial for A2 with small $\delta$ (range = 0.431), reflecting the high sensitivity to the specific files in each fold's test set.}
\label{tab:perfold_v7_nll}
\small
\begin{tabular}{@{}lcccccc@{}}
\toprule
Attack / Method & Fold 1 & Fold 2 & Fold 3 & Fold 4 & Fold 5 & Mean \\
\midrule
\multicolumn{7}{@{}l}{\textit{A1 Command Swap --- $S_{\text{mean}}$}} \\
 & 0.726 & 0.823 & 0.800 & 0.787 & 0.880 & 0.803 \\
\multicolumn{7}{@{}l}{\textit{A2 X $\delta$=0.001 --- $S_{\text{mean}}$}} \\
 & 0.984 & 0.653 & 0.604 & 0.688 & 0.554 & 0.697 \\
\multicolumn{7}{@{}l}{\textit{A2 X $\delta$=0.001 --- $S_{\max}$}} \\
 & 0.875 & 0.656 & 0.719 & 0.891 & 0.570 & 0.742 \\
\multicolumn{7}{@{}l}{\textit{A3 Parameter Swap --- $S_{\text{mean}}$}} \\
 & 0.642 & 0.553 & 0.603 & 0.648 & 0.658 & 0.621 \\
\multicolumn{7}{@{}l}{\textit{A4 Cross-Operation --- $S_{\text{mean}}$}} \\
 & 0.427 & 0.343 & 0.320 & 0.393 & 0.346 & 0.366 \\
\bottomrule
\end{tabular}
\end{table}

\begin{table}[ht]
\centering
\caption{V8 decoder reconstruction accuracy per fold. Fold-to-fold variation is modest for token accuracy (range = 3.1~pp) and numeric accuracy (range = 3.4~pp), but wider for command accuracy (range = 8.2~pp), reflecting sensitivity to the operation-class composition of each fold's training set.}
\label{tab:perfold_v8_training}
\begin{tabular}{@{}lccc@{}}
\toprule
Fold & Token Acc.\ (\%) & Numeric Acc.\ (\%) & Command Acc.\ (\%) \\
\midrule
1 & 65.6 & 42.6 & 84.1 \\
2 & 66.9 & 42.2 & 89.5 \\
3 & 66.7 & 43.1 & 90.7 \\
4 & 67.9 & 43.7 & 87.0 \\
5 & 68.7 & 45.6 & 92.3 \\
\midrule
\textbf{Mean} & \textbf{67.2} & \textbf{43.4} & \textbf{88.7} \\
\bottomrule
\end{tabular}
\end{table}

Fold-to-fold variation differs by attack type.
For A1 command swap ($S_{\text{mean}}$), AUROC ranges from 0.726 to 0.880 (range = 0.154), and the lowest fold (Fold~1) remains well above chance.
For A2 X $\delta = 0.001$ ($S_{\text{mean}}$), the range is much wider (0.554 to 0.984).
Fold~1 achieves near-perfect detection (0.984), while Fold~5 is near chance (0.554).
We attribute this variability to the composition of each fold's test set.
Files differ in their share of confident numeric tokens, the population with the highest per-token detectability.
With only $\sim$110 windows per fold, the confident/uncertain ratio can vary between folds.

\section{Vocabulary Structure}
\label{app:vocab}

The G-code vocabulary comprises 712 tokens organized into four semantic categories.
Table~\ref{tab:vocab_structure} provides the breakdown.

\begin{table}[ht]
\centering
\caption{Vocabulary composition (712 tokens total). The multi-head decoder structure assigns each token to one semantic category, enabling fine-grained anomaly attribution.}
\label{tab:vocab_structure}
\begin{tabular}{@{}lrl@{}}
\toprule
Category & Count & Examples \\
\midrule
Special tokens & 5 & PAD, BOS, EOS, UNK, SEP \\
Command tokens & $\sim$30 & G0, G1, G2, G3, M3, M5, M6, M30 \\
Parameter tokens & $\sim$10 & X, Y, Z, F, S, R, I, J, K \\
Numeric tokens & $\sim$667 & \texttt{NUM\_X\_1492}, \texttt{NUM\_Y\_0175}, \\
& & \texttt{NUM\_Z\_0600}, \texttt{NUM\_F\_2000} \\
\bottomrule
\end{tabular}
\end{table}

The multi-head output structure decouples these categories.
The type head selects among \{command, parameter, numeric, special\}, and the category-specific downstream heads predict within each category.
This factored structure reduces the effective vocabulary size per head; the command head, for example, predicts over $\sim$30 candidates rather than 712.
It also enables fine-grained anomaly attribution, because the per-head NLL breakdown indicates which token field is anomalous when the anomaly score is elevated.

\paragraph{Numeric tokenization.}
Numeric values (e.g., the coordinate X3.2750) are decomposed digit-by-digit using the \texttt{DigitByDigitValueHead}.
Each numeric token is decomposed into:
\begin{enumerate}
    \item A sign token ($+$ or $-$)
    \item Up to 4 integer digit tokens (hundreds, tens, ones, and a leading-zero position)
    \item A decimal point token
    \item Up to 4 decimal digit tokens (tenths, hundredths, thousandths, ten-thousandths)
\end{enumerate}
Each numeric token is represented by a sign position plus 6 digit positions, yielding 7 predictions per numeric value.
The digit-by-digit approach avoids the combinatorial explosion of discretizing continuous axis values into fixed bins.
Covering the machine's X-axis travel ($\approx$5.5~in) at the tokenizer's $0.001$~in grid would require $\sim$5{,}500 bins in a direct discretization scheme, compared to approximately 30 token types in the digit-by-digit scheme (10 digits + sign + decimal point, used combinatorially).

% Scoring method implementation details moved to supplementary material.

\section{Error Case Analysis}
\label{app:errors}

Three cases illustrate detection outcomes: a true positive, a false positive, and a false negative.

\paragraph{Illustrative true positive (A1 command swap).}
A sensor window during active face-milling contains the normal G-code sequence: G1 X-10.500 Y5.250 F200.
The decoder assigns $p = 0.72$ to G1 (the correct command) and $p = 0.01$ to G0 (the attacker's substitution).
The per-token NLL is 0.33 nats for the normal token and 4.61 nats for the attacked token, a $14\times$ increase.
The attacked window's $S_{\text{mean}}$ is 2.89 nats versus 1.47 nats for the normal window, well above typical detection thresholds.
This attack is highly detectable because the decoder has learned that G1 (feed-rate interpolation) is the expected motion command during active cutting and that G0 (rapid traverse) is anomalous in this context.

\paragraph{Illustrative false positive (tool-change sequence).}
A window during a tool-change sequence contains M6 (tool change) followed by G43 (tool-length compensation) and G0 (rapid positioning).
The M6 command occurs in fewer than 2\% of training windows, so the decoder assigns only $p \approx 0.15$ to M6 while assigning $p \approx 0.45$ to the more common G1.
This yields $S_{\text{mean}} = 4.2$ nats against a normal-window mean of 3.1 nats, triggering a false alarm at typical detection thresholds.
The false positive arises because rare but legitimate G-code patterns are underrepresented in training, not because the reconstruction is in error.
The decoder correctly identifies M6 as less probable than G1; the ranking is technically accurate but operationally unhelpful.

\paragraph{Illustrative false negative (A2 coordinate perturbation in uncertain regime).}
A coordinate perturbation of $\delta = 0.001$~in on an X-axis value (X3.2750 $\to$ X3.2760, one tokenizer grid step) falls in the decoder's low-confidence regime.
The decoder assigns $p = 0.003$ to the original value and $p = 0.002$ to the perturbed value, yielding $\Delta$NLL $= -\log(0.002) - (-\log(0.003)) = 0.41$ nats.
The original value ranks 47 out of 712 and the perturbed value 52 out of 712, a difference of only 5 ranks that yields $\Delta S_{\text{rank}} = 5/712 = 0.007$.
Both differences are well within the noise floor of the respective score distributions on normal data.
The V7 decoder cannot detect this window because its limited context (one G-code line) provides insufficient information to discriminate between X = 3.2750 and X = 3.2760, both of which are plausible axis values given the sensor observations.
The V8 decoder, by seeing the preceding and following G-code lines, may be able to detect this perturbation if the sequential pattern of X values is sufficiently constrained.

\section{Decoder Hyperparameters}
\label{app:hyperparameters}

Table~\ref{tab:hyperparameters} lists all decoder hyperparameters for both the V7 and V8 configurations.

\begin{table}[ht]
\centering
\caption{Decoder hyperparameters. Both V7 and V8 share the same architecture; they differ only in \texttt{max\_token\_len} and consequently in training epochs required for convergence.}
\label{tab:hyperparameters}
\small
\begin{tabular}{@{}llcc@{}}
\toprule
Category & Parameter & V7 & V8 \\
\midrule
\multirow{5}{*}{Architecture} & $d_{\text{model}}$ & 384 & 384 \\
& $n_{\text{layers}}$ & 8 & 8 \\
& $n_{\text{heads}}$ & 8 & 8 \\
& $d_{\text{ff}}$ & 1536 & 1536 \\
& \texttt{max\_token\_len} & 6 & 64 \\
\midrule
\multirow{6}{*}{Training} & Learning rate & $3 \times 10^{-4}$ & $3 \times 10^{-4}$ \\
& Weight decay & $1 \times 10^{-2}$ & $1 \times 10^{-2}$ \\
& Batch size & 32 & 32 \\
& Warmup epochs & 5 & 5 \\
& Total epochs & 50 & 80 \\
& Gradient clip norm & 1.0 & 1.0 \\
\midrule
\multirow{3}{*}{Regularization} & Dropout & 0.1 & 0.1 \\
& Label smoothing ($\epsilon$) & 0.1 & 0.1 \\
& Focal loss ($\gamma$) & 2.0 & 2.0 \\
\midrule
\multirow{3}{*}{Model size} & Decoder parameters & 22.1M & 22.1M \\
& Encoder parameters (frozen) & 5.1M & 5.1M \\
& Vocabulary size & 712 & 712 \\
\bottomrule
\end{tabular}
\end{table}

% V8 coordinate detail table moved to supplementary material.

The V8 $S_{\text{mean}}$ scores for coordinate attacks are generally lower than the V7 rank-based scores reported in the main text (Table~\ref{tab:main_results}), because the mean NLL is diluted by the many unperturbed tokens in the 63-token V8 sequence.
V8 $S_{\max}$ (not shown for all axes) achieves higher AUROC for larger $\delta$ values by isolating the single most anomalous token, as reported in the main results table.
